\documentclass[10pt,a4paper]{amsart}
\usepackage{amsmath,amssymb,amsthm,graphicx}
% Retained preamble for the Stacks Project chapter snapshot in this directory.
% The chapter's own preamble.tex is a copy of the shared Stacks Project preamble; the block
% below is sliced MECHANICALLY, line for line, from it. Nothing was added, removed or
% redefined; the wrapper class, the xr-hyper/hyperref wiring and the externaldocument list
% are omitted because the assembler supplies the document class.
%
% lines taken from preamble.tex: 3, 5, 9, 11, 12, 156, 157, 158, ... (51 lines total)
%
% package coverage restored from the chapter's own preamble (2026-10-13): xy, xy-2cell

\usepackage[all]{xy}
\xyoption{2cell}
\UseAllTwocells

\newenvironment{reference}{\comment}{\endcomment}
\newenvironment{slogan}{\comment}{\endcomment}
\newenvironment{history}{\comment}{\endcomment}
\theoremstyle{plain}
\newtheorem{theorem}[subsection]{Theorem}
\newtheorem{proposition}[subsection]{Proposition}
\newtheorem{lemma}[subsection]{Lemma}
\theoremstyle{definition}
\newtheorem{definition}[subsection]{Definition}
\newtheorem{example}[subsection]{Example}
\newtheorem{exercise}[subsection]{Exercise}
\newtheorem{situation}[subsection]{Situation}
\theoremstyle{remark}
\newtheorem{remark}[subsection]{Remark}
\newtheorem{remarks}[subsection]{Remarks}
\numberwithin{equation}{subsection}
\def\lim{\mathop{\mathrm{lim}}\nolimits}
\def\colim{\mathop{\mathrm{colim}}\nolimits}
\def\Spec{\mathop{\mathrm{Spec}}}
\def\Hom{\mathop{\mathrm{Hom}}\nolimits}
\def\Ext{\mathop{\mathrm{Ext}}\nolimits}
\def\SheafHom{\mathop{\mathcal{H}\!\mathit{om}}\nolimits}
\def\SheafExt{\mathop{\mathcal{E}\!\mathit{xt}}\nolimits}
\def\Sch{\mathit{Sch}}
\def\Mor{\mathop{\mathrm{Mor}}\nolimits}
\def\Ob{\mathop{\mathrm{Ob}}\nolimits}
\def\Sh{\mathop{\mathit{Sh}}\nolimits}
\def\NL{\mathop{N\!L}\nolimits}
\def\CH{\mathop{\mathrm{CH}}\nolimits}
\def\proetale{{pro\text{-}\acute{e}tale}}
\def\etale{{\acute{e}tale}}
\def\QCoh{\mathit{QCoh}}
\def\Ker{\mathop{\mathrm{Ker}}}
\def\Im{\mathop{\mathrm{Im}}}
\def\Coker{\mathop{\mathrm{Coker}}}
\def\Coim{\mathop{\mathrm{Coim}}}
\DeclareMathSymbol{\boxtimes}{\mathbin}{AMSa}{"02}
\def\QCohstack{\mathcal{QC}\!\mathit{oh}}
\def\Cohstack{\mathcal{C}\!\mathit{oh}}
\def\Spacesstack{\mathcal{S}\!\mathit{paces}}
\def\Quotfunctor{\mathrm{Quot}}
\def\Hilbfunctor{\mathrm{Hilb}}
\def\Curvesstack{\mathcal{C}\!\mathit{urves}}
\def\Polarizedstack{\mathcal{P}\!\mathit{olarized}}
\def\Complexesstack{\mathcal{C}\!\mathit{omplexes}}
\def\Pic{\mathop{\mathrm{Pic}}\nolimits}
\def\Picardstack{\mathcal{P}\!\mathit{ic}}
\def\Picardfunctor{\mathrm{Pic}}
\def\Deformationcategory{\mathcal{D}\!\mathit{ef}}
\usepackage{cleveref}
\crefname{originallemma}{original lemma}{original lemmas}
\Crefname{originallemma}{Original lemma}{Original lemmas}
\makeatletter
\newlabel{curves-lemma-characterize-contraction-to-stable}{{0E7Q}{source}{}{sourcepointer-1}{}}
\newlabel{curves-lemma-characterize-contraction-to-stable@cref}{{[originallemma][0][]0E7Q}{[1][0][]source}{}{}{}}
\makeatother
\crefname{originallemma}{original lemma}{original lemmas}
\Crefname{originallemma}{Original lemma}{Original lemmas}
\makeatletter
\newlabel{curves-section-contracting-to-stable}{{0E7N}{source}{}{sourcepointer-2}{}}
\newlabel{curves-section-contracting-to-stable@cref}{{[originallemma][0][]0E7N}{[1][0][]source}{}{}{}}
\makeatother
\crefname{originallemma}{original lemma}{original lemmas}
\Crefname{originallemma}{Original lemma}{Original lemmas}
\makeatletter
\newlabel{lemma-contract}{{0E88}{source}{}{sourcepointer-3}{}}
\newlabel{lemma-contract@cref}{{[originallemma][0][]0E88}{[1][0][]source}{}{}{}}
\makeatother
\crefname{originallemma}{original lemma}{original lemmas}
\Crefname{originallemma}{Original lemma}{Original lemmas}
\makeatletter
\newlabel{lemma-contract-basic}{{0E7C}{source}{}{sourcepointer-4}{}}
\newlabel{lemma-contract-basic@cref}{{[originallemma][0][]0E7C}{[1][0][]source}{}{}{}}
\makeatother
\crefname{originallemma}{original lemma}{original lemmas}
\Crefname{originallemma}{Original lemma}{Original lemmas}
\makeatletter
\newlabel{lemma-contract-basic-uniqueness}{{0E89}{source}{}{sourcepointer-5}{}}
\newlabel{lemma-contract-basic-uniqueness@cref}{{[originallemma][0][]0E89}{[1][0][]source}{}{}{}}
\makeatother
\crefname{originallemma}{original lemma}{original lemmas}
\Crefname{originallemma}{Original lemma}{Original lemmas}
\makeatletter
\newlabel{lemma-stable-curves}{{0E76}{source}{}{sourcepointer-6}{}}
\newlabel{lemma-stable-curves@cref}{{[originallemma][0][]0E76}{[1][0][]source}{}{}{}}
\makeatother
\crefname{originallemma}{original lemma}{original lemmas}
\Crefname{originallemma}{Original lemma}{Original lemmas}
\makeatletter
\newlabel{lemma-stable-one-piece-per-genus}{{0E78}{source}{}{sourcepointer-7}{}}
\newlabel{lemma-stable-one-piece-per-genus@cref}{{[originallemma][0][]0E78}{[1][0][]source}{}{}{}}
\makeatother
\title{Stacks, Moduli of Curves: lemma-contract-prestable-to-stable}
\author{The Stacks Project Authors}
\date{stacks-project:moduli-curves.tex:local-snapshot-20261012; GNU FDL 1.2 or later}
\begin{document}
\maketitle
\noindent\textit{Editorial note.} Statement span is lines 3126-3140 of sources/stacks-moduli-curves/pinned-original-source.tex; proof span is lines 3142-3192, the complete author proof (49 body lines).\par\medskip
\noindent\textit{Editorial note.} Navigation stubs: 7 cited-but-not-included labels are supplied as non-load-bearing pointers with their pinned tags where the snapshot tag index has them.\par\medskip
\noindent\textit{Editorial note.} extra\_packages is empty on purpose: the retained preamble already loads xy with options and a second declaration clashes in the assembler wrapper.\par\medskip
\noindent\textit{Editorial note.} Text is retained verbatim, including the author's own typos and placeholders.\par\medskip
\noindent\textit{Editorial note.} Queue index 6 of queue\_moduli-curves.json (pb=2307, stmt 3126-3140, proof 3142-3192).\par\medskip
\noindent\textit{Allowed standard prerequisites (editorial).}\begin{itemize}
\item Explicit prerequisite boundary. The cited inputs are prerequisites; the argument in the retained span is not replaced by any cited result.
\end{itemize}\medskip
\ifdefined\hypertarget\hypertarget{sourcepointer-1}{}\fi
\noindent\textit{Source reference.} Original lemma 0E7Q: sources/stacks-moduli-curves/pinned-original-source.tex, source lines 3192--3192. Cited by the retained statement or proof but not part of the collected spans; supplied only as a navigation pointer. The argument does not depend on its content.\par\medskip
\ifdefined\hypertarget\hypertarget{sourcepointer-2}{}\fi
\noindent\textit{Source reference.} Original lemma 0E7N: sources/stacks-moduli-curves/pinned-original-source.tex, source lines 3192--3192. Cited by the retained statement or proof but not part of the collected spans; supplied only as a navigation pointer. The argument does not depend on its content.\par\medskip
\ifdefined\hypertarget\hypertarget{sourcepointer-3}{}\fi
\noindent\textit{Source reference.} Original lemma 0E88: sources/stacks-moduli-curves/pinned-original-source.tex, source lines 2865--2874. Cited by the retained statement or proof but not part of the collected spans; supplied only as a navigation pointer. The argument does not depend on its content.\par\medskip
\ifdefined\hypertarget\hypertarget{sourcepointer-4}{}\fi
\noindent\textit{Source reference.} Original lemma 0E7C: sources/stacks-moduli-curves/pinned-original-source.tex, source lines 2987--2995. Cited by the retained statement or proof but not part of the collected spans; supplied only as a navigation pointer. The argument does not depend on its content.\par\medskip
\ifdefined\hypertarget\hypertarget{sourcepointer-5}{}\fi
\noindent\textit{Source reference.} Original lemma 0E89: sources/stacks-moduli-curves/pinned-original-source.tex, source lines 2900--2912. Cited by the retained statement or proof but not part of the collected spans; supplied only as a navigation pointer. The argument does not depend on its content.\par\medskip
\ifdefined\hypertarget\hypertarget{sourcepointer-6}{}\fi
\noindent\textit{Source reference.} Original lemma 0E76: sources/stacks-moduli-curves/pinned-original-source.tex, source lines 2679--2688. Cited by the retained statement or proof but not part of the collected spans; supplied only as a navigation pointer. The argument does not depend on its content.\par\medskip
\ifdefined\hypertarget\hypertarget{sourcepointer-7}{}\fi
\noindent\textit{Source reference.} Original lemma 0E78: sources/stacks-moduli-curves/pinned-original-source.tex, source lines 2736--2749. Cited by the retained statement or proof but not part of the collected spans; supplied only as a navigation pointer. The argument does not depend on its content.\par\medskip
\begin{lemma}
\label{lemma-contract-prestable-to-stable}
Let $f : X \to S$ be a prestable family of curves of genus $g \geq 2$.
There is a factorization $X \to Y \to S$ of $f$ where $g : Y \to S$
is a stable family of curves and $c : X \to Y$ has the following
properties
\begin{enumerate}
\item $\mathcal{O}_Y = c_*\mathcal{O}_X$ and $R^1c_*\mathcal{O}_X = 0$
and this remains true after base change by any morphism $S' \to S$, and
\item for any $s \in S$ the morphism $c_s : X_s \to Y_s$ is the
contraction of rational tails and bridges discussed in
Algebraic Curves, Section \ref{curves-section-contracting-to-stable}.
\end{enumerate}
Moreover $c : X \to Y$ is unique up to unique isomorphism.
\end{lemma}

\begin{proof}
Let $s \in S$. Let $c_0 : X_s \to Y_0$ be the contraction of
Algebraic Curves, Section \ref{curves-section-contracting-to-stable}
(more precisely Algebraic Curves, Lemma
\ref{curves-lemma-characterize-contraction-to-stable}).
By Lemma \ref{lemma-contract-basic}
there exists an elementary \'etale neighbourhood
$(U, u)$ and a morphism $c : X_U \to Y$
of families of curves over $U$ which recovers
$c_0$ as the fibre at $u$.
Since $\omega_{Y_0}$ is ample, after possibly shrinking $U$,
we see that $Y \to U$ is a stable family of genus $g$
by the openness inherent in
Lemmas \ref{lemma-stable-curves} and \ref{lemma-stable-one-piece-per-genus}.
After possibly shrinking $U$ once more, assertion (1) of the lemma for
$c : X_U \to Y$ follows from Lemma \ref{lemma-contract}.
Moreover, part (2) holds by the uniqueness in Algebraic Curves, Lemma
\ref{curves-lemma-characterize-contraction-to-stable}.
We conclude that a morphism $c$ as in the lemma exists \'etale locally
on $S$. More precisely, there exists an \'etale covering
$\{U_i \to S\}$ and morphisms $c_i : X_{U_i} \to Y_i$ over $U_i$
where $Y_i \to U_i$ is a stable family of curves
having properties (1) and (2) stated in the lemma.

\medskip\noindent
To finish the proof it suffices to prove uniqueness of $c : X \to Y$
(up to unique isomorphism). Namely, once this is done, then we
obtain isomorphisms
$$
\varphi_{ij} :
Y_i \times_{U_i} (U_i \times_S U_j)
\longrightarrow
Y_i \times_{U_j} (U_i \times_S U_j)
$$
satisfying the cocycle condition (by uniqueness) over
$U_i \times U_j \times U_k$. Since $\overline{\mathcal{M}_g}$
is an algebraic stack, we have effectiveness of descent data
and we obtain $Y \to S$. The morphisms $c_i$ descend to a morphism
$c : X \to Y$ over $S$. Finally, properties (1) and (2) for $c$
are immediate from properties (1) and (2) for $c_i$.

\medskip\noindent
Finally, if $c_1 : X \to Y_i$, $i = 1, 2$ are two morphisms towards
stably families of curves over $S$ satisfying (1) and (2), then
we obtain a morphism $Y_1 \to Y_2$ compatible with $c_1$ and $c_2$
at least locally on $S$ by Lemma \ref{lemma-contract-basic-uniqueness}.
We omit the verification that these morphisms are unique
(hint: this follows from the fact that the scheme theoretic image
of $c_1$ is $Y_1$). Hence these locally given morphisms glue
and the proof is complete.
\end{proof}

\end{document}
